\documentclass[11pt,letterpaper]{article}

% Essential Packages
\usepackage[margin=1in]{geometry}
\usepackage{amsmath,amssymb}
\usepackage{graphicx}
\usepackage{circuitikz}
\usepackage{tikz}
\usetikzlibrary{shapes,arrows,positioning,calc}
\usepackage{booktabs}
\usepackage{multirow}
\usepackage{xcolor}
\usepackage{fancyhdr}
\usepackage{tcolorbox}
\usepackage{enumitem}
\usepackage{float}
\usepackage{listings}
\usepackage{hyperref}

% Colors
\definecolor{nlcblue}{RGB}{0,51,102}
\definecolor{alertred}{RGB}{204,0,0}
\definecolor{safgreen}{RGB}{0,128,0}

% Header/Footer
\pagestyle{fancy}
\fancyhead[L]{\textbf{NLC STEM Club UAV Project}}
\fancyhead[R]{\textbf{Technical Terminology \& Definitions}}
\fancyfoot[C]{\thepage}

% Title
\title{\Huge\textbf{Technical Terminology \& Definitions}\\
\Large Reference Guide for Electrical Team\\
\large NLC STEM Club - Electrical Systems Team}
\author{February 2026}
\date{}

\begin{document}

\maketitle
\thispagestyle{empty}
\newpage
\section{TECHNICAL TERMINOLOGY \& DEFINITIONS}

\begin{itemize}
  \item Fixed-Wing Environmental Monitoring UAV
  \item NLC STEM Club - Electrical Systems Team
  \item Reference Guide for Team Members
\end{itemize}

\section{DOCUMENT PURPOSE:}

\begin{itemize}
  \item This glossary defines all technical terms, acronyms, and concepts used throughout
  \item the UAV electrical system documentation. Use this as a quick reference guide when
  \item encountering unfamiliar terminology.
\end{itemize}

\section{ORGANIZATION:}

\begin{itemize}
  \item Terms are organized alphabetically within categories for easy lookup.
\end{itemize}

\section{SECTION 1: ELECTRICAL TERMS \& CONCEPTS}

AMPERE (A) / MILLIAMPERE (mA)\par

\begin{itemize}
  \item Unit of electrical current. 1 Ampere = 1000 milliamperes.
  \item Example: The flight controller draws 250mA (0.25A) of current.
\end{itemize}

AWG (American Wire Gauge)\par

\begin{itemize}
  \item Standard system for measuring wire thickness. LOWER numbers = THICKER wire.
  \item 12 AWG: Very thick, high current (30A+) - motor wires
  \item 18 AWG: Medium, moderate current (5-10A) - power distribution
  \item 26 AWG: Thin, low current/signal (<1A) - sensor wires
  \item The thicker the wire, the less resistance and more current it can safely carry.
\end{itemize}

\section{BACKFEEDING}

\begin{itemize}
  \item Unwanted flow of power in the reverse direction through a circuit. In our system,
  \item this occurs if the ESC's +5V wire is connected to the FC, potentially damaging
  \item the FC's voltage regulator. SOLUTION: Cut the middle wire on ESC signal cable.
\end{itemize}

BEC (Battery Eliminator Circuit)\par

\begin{itemize}
  \item A voltage regulator that converts high battery voltage (14.8V) down to 5V for
  \item electronics. "Eliminates" the need for a separate battery for radio equipment.
  \item Our ESC has a built-in 5V 3A BEC (Linear BEC - less efficient but simpler).
\end{itemize}

\section{BROWNOUT}

\begin{itemize}
  \item Temporary voltage drop that causes electronics to malfunction or reset. Occurs
  \item when too much current is drawn from a power source, causing voltage to sag below
  \item minimum operating level. Our Dual-BEC design prevents FC brownout from servo loads.
\end{itemize}

CAPACITY (mAh or Ah)\par

\begin{itemize}
  \item Milliamp-hours or Amp-hours. Measures how much charge a battery can store.
  \item Our battery: 2200mAh = can deliver 2200mA for 1 hour (or 1100mA for 2 hours)
  \item Higher capacity = longer flight time (but also heavier battery)
\end{itemize}

\section{CONTINUITY}

\begin{itemize}
  \item Unbroken electrical path. Used in testing: "check continuity" means verify a
  \item wire is not broken and conducts electricity. Measured with multimeter in $\Omega$ mode.
\end{itemize}

\section{CURRENT (I)}

\begin{itemize}
  \item Flow of electrical charge, measured in Amperes (A). Like water flow in a pipe.
  \item Formula: I = V / R (Current = Voltage / Resistance)
  \item Typical current: FC draws 250mA
  \item Peak current: Motor can draw 30A during full throttle
\end{itemize}

\section{DUAL-BEC CONFIGURATION}

\subsection{Our chosen power architecture using TWO separate voltage regulators}

\begin{itemize}
  \item BEC 1 (ESC integrated): Powers flight-critical systems (FC, GPS, RX)
  \item BEC 2 (External UBEC): Powers high-current servos
  \item Advantage: Isolates servo power spikes from flight controller, preventing crashes.
\end{itemize}

\section{DUTY CYCLE}

\begin{itemize}
  \item Percentage of time a signal is "ON" vs "OFF". Used in PWM control.
  \item 50\% duty cycle = signal is high 50\% of the time
  \item Servo control: 1ms pulse (5\% duty cycle) = full left, 2ms (10\%) = full right
\end{itemize}

ELECTRICAL NOISE / EMI (Electromagnetic Interference)\par

\subsection{Unwanted electrical signals that interfere with sensitive electronics. Sources}

\begin{itemize}
  \item Motor/ESC switching (high-frequency noise)
  \item Power supply ripple
  \item Radio frequency interference
  \item Mitigation: Twisted wires, distance separation, filtering capacitors.
\end{itemize}

GND (Ground)\par

\begin{itemize}
  \item Reference point for voltage measurement, typically 0V. Also completes circuit
  \item for current flow. CRITICAL: All components must share a common ground or signals
  \item will be corrupted. Our design uses star grounding topology.
\end{itemize}

LOAD\par

\subsection{Any device that consumes electrical power. In our system}

\begin{itemize}
  \item Light load: GPS (60mA)
  \item Medium load: Arduino (263mA)
  \item Heavy load: Motor (10-30A)
\end{itemize}

\section{OHMS LAW}

\begin{itemize}
  \item Fundamental relationship: V = I $\times$ R (Voltage = Current $\times$ Resistance)
  \item Also: I = V / R and R = V / I
  \item Used to calculate current draw, voltage drop, wire sizing, etc.
\end{itemize}

\section{PARALLEL CONNECTION}

\begin{itemize}
  \item Components connected with all positive terminals together and all negative
  \item terminals together. Voltage stays same, currents add.
  \item Example: All 4 servos powered in parallel from UBEC (each gets 5V, currents sum).
\end{itemize}

\section{POLARITY}

\begin{itemize}
  \item Direction of electrical flow. Positive (+) and Negative (-) terminals.
  \item [WARNING] CRITICAL: Reversed polarity can instantly destroy electronics!
  \item Always verify RED = positive, BLACK = negative before connecting.
\end{itemize}

POWER (W, Watts)\par

\begin{itemize}
  \item Rate of energy consumption. Formula: P = V $\times$ I (Power = Voltage $\times$ Current)
  \item Example: Motor at 14.8V drawing 20A = 296W power consumption
  \item Our battery: 2200mAh $\times$ 14.8V = 32.56Wh (Watt-hours) total energy
\end{itemize}

\section{POWER RAIL}

\begin{itemize}
  \item Distribution bus that supplies power to multiple components at same voltage.
\end{itemize}

\subsection{Our system has two 5V rails}

\begin{itemize}
  \item Rail 1: Flight-critical (FC, GPS, RX, Arduino) - ESC BEC
  \item Rail 2: Servo power (all 4 servos) - External UBEC
\end{itemize}

\section{RESISTANCE (R)}

\begin{itemize}
  \item Opposition to current flow, measured in Ohms ($\Omega$). Like friction in a pipe.
  \item Low resistance: Thick wires, good connections
  \item High resistance: Thin wires, corroded connections
  \item Causes voltage drop and heat generation.
\end{itemize}

\section{SERIES CONNECTION}

\begin{itemize}
  \item Components connected end-to-end so current flows through each in sequence.
  \item Voltages add, current stays same.
  \item Example: Battery cells in series (3.7V + 3.7V + 3.7V + 3.7V = 14.8V for 4S).
\end{itemize}

\section{SHORT CIRCUIT}

\begin{itemize}
  \item Unintended direct connection between positive and negative, bypassing load.
  \item Results in VERY high current, heat, fire risk, component damage.
  \item [WARNING] ALWAYS use fuse/protection and inspect for bare wire touching!
\end{itemize}

\section{STAR GROUNDING}

\begin{itemize}
  \item Grounding topology where all ground wires connect to ONE central point (battery
  \item negative) rather than daisy-chaining. Prevents ground loops and voltage differences
  \item between components. REQUIRED for our system's proper operation.
\end{itemize}

UBEC (Universal Battery Eliminator Circuit)\par

\begin{itemize}
  \item Advanced switching voltage regulator (vs linear BEC). More efficient, handles
  \item higher input voltage range, generates less heat. Our external UBEC is 5V 6A
  \item switching regulator for servo power.
\end{itemize}

\section{VOLTAGE (V)}

\begin{itemize}
  \item Electrical potential difference, measured in Volts. Like water pressure in a pipe.
  \item Battery voltage: 14.8V nominal (13.2-16.8V range for 4S LiPo)
  \item Regulated voltage: 5.0V for electronics
  \item Signal voltage: 3.3V or 5V logic levels
\end{itemize}

\section{VOLTAGE DROP}

\begin{itemize}
  \item Reduction in voltage due to resistance in wires/connections. Formula: V\_drop = I $\times$ R
  \item Problem: Long thin wires carrying high current have significant voltage drop.
  \item Solution: Use appropriately thick wire for current level (AWG sizing).
\end{itemize}

\section{VOLTAGE SAG}

\begin{itemize}
  \item Temporary voltage drop when high current is suddenly drawn. Battery internal
  \item resistance causes voltage to "sag" under heavy load. Can cause brownout if
  \item excessive. Our high C-rating battery (130C) minimizes voltage sag.
\end{itemize}

\section{SECTION 2: COMPONENT TERMINOLOGY}

\section{ARDUINO}

\begin{itemize}
  \item Open-source microcontroller platform. Our Arduino Nano uses ATmega328P chip
  \item (8-bit, 16MHz, 32KB flash) for data logging and sensor interfacing. Programmed
  \item in C/C++ using Arduino IDE.
\end{itemize}

\section{BAROMETER}

\begin{itemize}
  \item Atmospheric pressure sensor (BMP280 on flight controller). Measures air pressure
  \item to calculate altitude. Used for altitude hold and vertical navigation.
\end{itemize}

BLDC MOTOR (Brushless DC Motor)\par

\begin{itemize}
  \item Electric motor with electronic commutation (no brushes). More efficient, reliable,
  \item and powerful than brushed motors. Requires ESC for speed control. Our D2830EVO
  \item is a 3-phase outrunner BLDC motor.
\end{itemize}

C-RATING (Battery)\par

\begin{itemize}
  \item Discharge rate capability of LiPo battery. Formula: Max Current = Capacity $\times$ C-Rating
  \item Our battery: 2.2Ah $\times$ 130C = 286A maximum burst current (never approach this!)
  \item Higher C-rating = lower internal resistance = less voltage sag under load.
\end{itemize}

COMPASS (Magnetometer)\par

\begin{itemize}
  \item Sensor that detects Earth's magnetic field for heading determination. Integrated
  \item in BN-220 GPS module (HMC5883L chip). Requires calibration and distance from
  \item magnetic interference (motors, current-carrying wires).
\end{itemize}

ESC (Electronic Speed Controller)\par

\begin{itemize}
  \item Power electronics that control brushless motor speed. Converts DC battery power
  \item to 3-phase AC for motor. Receives PWM throttle signal from FC, outputs variable
  \item frequency/amplitude to motor. Our 40A ESC has integrated 5V BEC.
\end{itemize}

FC (Flight Controller)\par

\begin{itemize}
  \item Central computer that stabilizes aircraft and executes flight commands. Our
  \item MATEK F405-WMN has:
  \item STM32F405 processor (32-bit ARM, 168MHz)
  \item IMU (gyro + accelerometer) for attitude sensing
  \item Runs iNav or ArduPilot firmware
\end{itemize}

\section{FOLDING PROPELLER}

\begin{itemize}
  \item Propeller with blades that fold back when motor is off (reduces drag in glide).
  \item Spring-loaded to deploy under centrifugal force when spinning. Used on our
  \item motor for efficient glide performance.
\end{itemize}

GPS (Global Positioning System)\par

\begin{itemize}
  \item Satellite navigation system. Our BN-220 module receives signals from GPS
  \item satellites to determine 3D position (latitude, longitude, altitude) and velocity.
  \item Requires clear view of sky; 6+ satellites for good fix.
\end{itemize}

\section{GYROSCOPE}

\begin{itemize}
  \item Sensor that measures angular rotation rate (degrees per second). Part of IMU
  \item on flight controller. Used for stabilization - detects when aircraft is rotating
  \item and commands control surfaces to counteract.
\end{itemize}

\section{HEAT SHRINK TUBING}

\begin{itemize}
  \item Plastic tubing that shrinks when heated, used to insulate electrical connections.
  \item Always cover solder joints and bare wire with heat shrink to prevent shorts.
\end{itemize}

IMU (Inertial Measurement Unit)\par

\begin{itemize}
  \item Combined sensor package: 3-axis gyroscope + 3-axis accelerometer = 6 DOF
  \item (degrees of freedom). Measures rotation and acceleration in all directions.
  \item Flight controller uses IMU data for stabilization and attitude estimation.
\end{itemize}

\section{JST CONNECTOR}

\subsection{Small standardized connector series. Common types}

\begin{itemize}
  \item JST-XH: Battery balance connector
  \item JST-SH: 1.0mm pitch (GPS, receiver)
  \item JST-PH: 2.0mm pitch (general electronics)
\end{itemize}

KV RATING (Motor)\par

\begin{itemize}
  \item RPM per volt constant. Our 1000KV motor spins at 1000 RPM for every volt applied.
  \item At 14.8V: 14,800 RPM (no load)
  \item Lower KV = more torque, larger props; Higher KV = less torque, smaller props.
\end{itemize}

LiPo (Lithium Polymer Battery)\par

\begin{itemize}
  \item Rechargeable battery with high energy density. Chemistry: Lithium-ion with
  \item polymer electrolyte. Nominal voltage: 3.7V per cell.
  \item 1S = 3.7V, 2S = 7.4V, 3S = 11.1V, 4S = 14.8V (our battery)
  \item [WARNING] DANGEROUS: Can catch fire if damaged, overcharged, or short-circuited!
\end{itemize}

OSD (On-Screen Display)\par

\begin{itemize}
  \item Overlays flight telemetry on FPV video feed. MATEK F405-WMN has integrated OSD
  \item chip that can display battery voltage, current, altitude, GPS data, etc. on
  \item screen. Configured in iNav/ArduPilot.
\end{itemize}

\section{OUTRUNNER MOTOR}

\begin{itemize}
  \item BLDC motor where outer case rotates and center is stationary (opposite of
  \item inrunner). Higher torque, better for direct-drive propellers. Our D2830EVO is
  \item an outrunner.
\end{itemize}

\section{RECEIVER (RX)}

\begin{itemize}
  \item Radio receiver that picks up transmitter signals. Our HAPPYMODEL EP2 receives
  \item ExpressLRS 2.4GHz radio commands and sends them to FC via CRSF or SBUS protocol.
\end{itemize}

\section{SENSOR FUSION}

\begin{itemize}
  \item Algorithm that combines data from multiple sensors (GPS, IMU, barometer, compass)
  \item to get better position/attitude estimate than any single sensor. Flight controller
  \item firmware does this automatically.
\end{itemize}

SERVO\par

\begin{itemize}
  \item Electromechanical actuator with position feedback. Receives PWM pulse (1-2ms)
  \item and rotates output shaft to corresponding angle. Our ES08MA II servos have:
  \item Metal gears (durable)
  \item 180$^\circ$ rotation
  \item Proportional control
\end{itemize}

\section{TRANSMITTER (TX)}

\begin{itemize}
  \item Hand-held radio controller. Pilot inputs (sticks, switches) converted to radio
  \item signals and sent to receiver on aircraft. Can also refer to data transmit pin
  \item on serial connection.
\end{itemize}

UBEC (Previously Defined)\par

\begin{itemize}
  \item See Section 1. External voltage regulator providing 5V @ 6A for servo power.
\end{itemize}

\section{XT60 CONNECTOR}

\begin{itemize}
  \item High-current connector rated for 60A continuous. Yellow plastic housing with
  \item spring-pin contacts. Standard for 3S-6S LiPo batteries. Gender: Male on battery,
  \item Female on ESC/cables.
\end{itemize}

\section{SECTION 3: PROTOCOLS \& COMMUNICATION STANDARDS}

\section{ANALOG SIGNAL}

\begin{itemize}
  \item Continuously variable voltage representing data. Our MQ135 gas sensor outputs
  \item 0-5V analog voltage proportional to gas concentration. Read by Arduino ADC
  \item (Analog-to-Digital Converter).
\end{itemize}

\section{BAUD RATE}

\subsection{Speed of serial communication in bits per second (bps). Common rates}

\begin{itemize}
  \item 9600 baud: GPS default (slow but reliable)
  \item 115200 baud: Fast serial (firmware upload, debugging)
  \item 420000 baud: CRSF protocol from ELRS receiver
  \item Both devices must use SAME baud rate to communicate.
\end{itemize}

CRSF (Crossfire Protocol)\par

\begin{itemize}
  \item Bi-directional serial protocol for RC control and telemetry. Developed by TBS,
  \item now used by ExpressLRS. 420000 baud, low latency, supports up to 16 channels
  \item plus telemetry data. Our ELRS receiver uses CRSF to talk to FC.
\end{itemize}

\section{DIGITAL SIGNAL}

\begin{itemize}
  \item Discrete on/off voltage levels representing binary data (0 or 1).
  \item Logic LOW: 0V (binary 0)
  \item Logic HIGH: 3.3V or 5V (binary 1)
  \item Used for: I2C, SPI, UART, PWM control signals.
\end{itemize}

ELRS (ExpressLRS)\par

\subsection{Open-source long-range radio control system using 2.4GHz or 900MHz. Features}

\begin{itemize}
  \item Very low latency (5-20ms)
  \item Long range (10+ km)
  \item Frequency hopping for interference immunity
  \item Over-the-air updates
  \item Our EP2 receiver uses ELRS protocol.
\end{itemize}

I2C (Inter-Integrated Circuit)\par

\begin{itemize}
  \item Two-wire serial protocol for short-distance communication. Uses SDA (data) and
  \item SCL (clock) lines. Multiple devices on same bus (each has unique address).
  \item Our BME280 sensor uses I2C to send temperature/humidity/pressure data to Arduino.
\end{itemize}

NMEA (National Marine Electronics Association)\par

\subsection{Text-based protocol for GPS data. Human-readable sentences like}

\section{\$GPGGA,123519,4807.038,N,01131.000,E,1,08,0.9,545.4,M,46.9,M,,*47}

\begin{itemize}
  \item Contains position, time, satellites, etc. Alternative to binary UBX protocol.
\end{itemize}

PWM (Pulse Width Modulation)\par

\subsection{Signal where information is encoded in pulse duration. Used for}

\begin{itemize}
  \item Servo control: 1ms = -90$^\circ$, 1.5ms = center, 2ms = +90$^\circ$
  \item Motor throttle: 1ms = 0\%, 2ms = 100\%
  \item Frequency: typically 50Hz (20ms period) for servos
  \item Analog value represented by digital signal duty cycle.
\end{itemize}

SBUS (Serial Bus)\par

\begin{itemize}
  \item Inverted serial protocol from FrSky. 100000 baud, 25ms update, 16 channels.
  \item Less efficient than CRSF but widely compatible. Fallback option if CRSF doesn't
  \item work with our FC firmware.
\end{itemize}

SPI (Serial Peripheral Interface)\par

\begin{itemize}
  \item Four-wire serial protocol: MOSI (master out), MISO (master in), SCK (clock),
  \item CS (chip select). Fast, full-duplex communication. Our MicroSD card module uses
  \item SPI to transfer data at high speed from Arduino.
\end{itemize}

\section{TELEMETRY}

\subsection{Real-time transmission of aircraft data back to ground station. Can include}

\begin{itemize}
  \item battery voltage, current, altitude, GPS position, speed, RSSI (signal strength).
  \item CRSF protocol supports bi-directional telemetry.
\end{itemize}

UART (Universal Asynchronous Receiver-Transmitter)\par

\begin{itemize}
  \item Serial communication protocol using TX (transmit) and RX (receive) wires.
  \item Asynchronous = no clock signal, must agree on baud rate.
  \item UART1 on FC: Connected to GPS (9600 baud)
  \item UART2 on FC: Connected to receiver (420000 baud)
\end{itemize}

\section{UBX PROTOCOL}

\begin{itemize}
  \item Binary GPS protocol from u-blox. More efficient than NMEA, higher update rate,
  \item configurable messages. Our BN-220 GPS supports both UBX and NMEA.
\end{itemize}

\section{SECTION 4: AVIATION \& UAV TERMINOLOGY}

\section{AILERON}

\begin{itemize}
  \item Control surface on wing trailing edge that controls roll (banking). Differential
  \item movement: left aileron up + right aileron down = roll left.
  \item Our aircraft has 2 ailerons (Servo 1 and Servo 2).
\end{itemize}

\section{AUTONOMOUS FLIGHT}

\begin{itemize}
  \item Aircraft flies pre-programmed mission without pilot input. Uses GPS waypoints,
  \item altitude commands, loiter circles. FC firmware executes mission autonomously.
  \item Requires GPS lock and properly tuned flight controller.
\end{itemize}

\section{BLACKBOX}

\begin{itemize}
  \item Flight data recorder that logs sensor data at high rate (1000Hz+) for post-
  \item flight analysis. Can be onboard flash memory or microSD card. Used for tuning
  \item PIDs, diagnosing crashes, analyzing performance.
\end{itemize}

\section{CONTROL SURFACE}

\subsection{Movable aerodynamic surface that controls aircraft attitude}

\begin{itemize}
  \item Ailerons: Roll control
  \item Elevator: Pitch control
  \item Rudder: Yaw control
\end{itemize}

\section{ELEVATOR}

\begin{itemize}
  \item Control surface on horizontal tail that controls pitch (nose up/down).
  \item Our Servo 3 controls the elevator.
\end{itemize}

ELRS (Already Defined)\par

\begin{itemize}
  \item See Section 3 - ExpressLRS long-range control system.
\end{itemize}

\section{FAILSAFE}

\subsection{Pre-programmed behavior when radio link is lost. Options}

\begin{itemize}
  \item Hold position (loiter)
  \item Return to home (RTH) and land
  \item Immediate land
  \item Cut throttle (glide down)
  \item Configured in flight controller firmware.
\end{itemize}

\section{FIXED-WING}

\begin{itemize}
  \item Aircraft with continuous rigid wings (like airplane), as opposed to rotorcraft
  \item (quadcopter). Our UAV is fixed-wing design. Requires forward motion for lift;
  \item can glide efficiently but cannot hover.
\end{itemize}

FPV (First Person View)\par

\begin{itemize}
  \item Flying from onboard camera perspective. Video transmitted to pilot's goggles/
  \item screen in real-time. Our project focuses on autonomous flight, not FPV, but
  \item FC has OSD capability for future FPV addition.
\end{itemize}

\section{GPS LOCK / FIX}

\begin{itemize}
  \item Successful acquisition of GPS satellite signals allowing position determination.
\end{itemize}

\subsection{Requires}

\begin{itemize}
  \item Clear sky view
  \item 4+ satellites for 2D fix (lat/long)
  \item 6+ satellites for 3D fix (lat/long/altitude)
  \item Indicated by LED on GPS module.
\end{itemize}

IMU (Already Defined)\par

\begin{itemize}
  \item See Section 2 - Inertial Measurement Unit for stabilization.
\end{itemize}

\section{LOITER / POSITION HOLD}

\begin{itemize}
  \item Flight mode where aircraft maintains fixed GPS position (circles within radius).
  \item Requires GPS lock. Useful for observing area or waiting for telemetry.
\end{itemize}

PID (Proportional-Integral-Derivative)\par

\begin{itemize}
  \item Control algorithm used for stabilization. FC continuously adjusts control surfaces
  \item to maintain desired attitude:
  \item P: Correction proportional to error (immediate response)
  \item I: Correction based on accumulated error (eliminates steady-state error)
  \item D: Correction based on rate of error change (dampening)
  \item Tuning PID values is critical for stable flight.
\end{itemize}

PITCH\par

\begin{itemize}
  \item Rotation about lateral axis (nose up/down). Controlled by elevator.
\end{itemize}

\section{PROPELLER}

\subsection{Rotating blade that converts motor power to thrust. Specifications}

\begin{itemize}
  \item Diameter $\times$ Pitch: 1060 = 10" diameter, 6" pitch
  \item Pitch = distance propeller advances in one rotation
  \item Folding prop: Blades fold back when not spinning
\end{itemize}

ROLL\par

\begin{itemize}
  \item Rotation about longitudinal axis (banking left/right). Controlled by ailerons.
\end{itemize}

RTH (Return to Home)\par

\begin{itemize}
  \item Autonomous flight mode where aircraft flies back to takeoff location using GPS.
  \item Can be triggered by pilot command, failsafe, or low battery.
\end{itemize}

\section{RUDDER}

\begin{itemize}
  \item Control surface on vertical tail that controls yaw (nose left/right).
  \item Our Servo 4 controls the rudder. Less critical in calm wind than ailerons/elevator.
\end{itemize}

\section{STABILIZATION}

\begin{itemize}
  \item Automatic correction to maintain level flight. Gyros detect rotation, FC commands
  \item control surfaces to counteract. Modes:
  \item Acro: No stabilization (manual input only)
  \item Angle: Self-leveling (returns to level when sticks centered)
  \item Horizon: mix of Acro and Angle
\end{itemize}

UAV (Unmanned Aerial Vehicle)\par

\begin{itemize}
  \item Aircraft without human pilot onboard. Controlled remotely or autonomously.
  \item Synonyms: drone, UAS (Unmanned Aircraft System), remotely piloted aircraft.
\end{itemize}

YAW\par

\begin{itemize}
  \item Rotation about vertical axis (nose left/right). Controlled by rudder.
\end{itemize}

\section{SECTION 5: SENSOR \& MEASUREMENT TERMINOLOGY}

ADC (Analog-to-Digital Converter)\par

\begin{itemize}
  \item Circuit that converts analog voltage to digital number. Arduino Nano has 10-bit
  \item ADC (0-1023 range) for reading analog sensors. MQ135 gas sensor output is read
  \item by ADC on pin A0.
\end{itemize}

\section{BME280}

\subsection{Combined environmental sensor IC (Bosch Sensortec)}

\begin{itemize}
  \item Barometric pressure: 300-1100 hPa (altitude 0-30,000 ft)
  \item Temperature: -40 to +85$^\circ$C ($\pm$1$^\circ$C accuracy)
  \item Humidity: 0-100\% RH ($\pm$3\% accuracy)
  \item Uses I2C interface. Perfect for weather monitoring mission.
\end{itemize}

\section{BURN-IN PERIOD}

\begin{itemize}
  \item Initial operation time required for sensor to stabilize. MQ135 gas sensor requires
  \item 24-48 hours of continuous power before accurate readings. Tin oxide sensor
  \item element needs to heat up and stabilize chemically.
\end{itemize}

\section{CALIBRATION}

\subsection{Procedure to adjust sensor readings to known reference values. Examples}

\begin{itemize}
  \item Compass: Rotate aircraft through 360$^\circ$ to map magnetic field distortions
  \item ESC: Set throttle range endpoints (min/max PWM values)
  \item Gas sensor: Expose to fresh air and record baseline reading
\end{itemize}

CSV (Comma-Separated Values)\par

\begin{itemize}
  \item Simple text file format for data logging. Each line is one record, values
  \item separated by commas. Example:
  \item Time,Temp,Humidity,AirQuality
  \item 0,24.5,65,450
  \item Our Arduino saves sensor data in CSV format for analysis in Excel/Python.
\end{itemize}

\section{DATA LOGGING}

\subsection{Recording sensor measurements over time for post-flight analysis. Our system}

\begin{itemize}
  \item Timestamp from millis() function
  \item BME280 readings every second
  \item MQ135 readings every second
  \item Written to microSD card in CSV format
\end{itemize}

\section{I2C ADDRESS}

\subsection{Unique 7-bit number identifying device on I2C bus. BME280 default address}

\begin{itemize}
  \item 0x76 or 0x77 (solder jumper selectable). Scanner sketch can find device addresses.
\end{itemize}

MicroSD CARD\par

\begin{itemize}
  \item Small removable flash storage. Our module uses SPI interface for data transfer.
  \item Format: FAT32 recommended (up to 32GB cards). Stores sensor data logs as CSV files.
\end{itemize}

MQ135\par

\subsection{Metal oxide gas sensor for air quality monitoring. Detects}

\begin{itemize}
  \item CO2 (carbon dioxide)
  \item NH3 (ammonia)
  \item NOx (nitrogen oxides)
  \item Alcohol, smoke, benzene
  \item Outputs analog voltage (0-5V) proportional to gas concentration. Requires 5V
  \item supply and 24-48 hour burn-in period. Not highly selective (detects multiple gases).
\end{itemize}

\section{RESOLUTION}

\subsection{Smallest change a sensor can detect. Examples}

\begin{itemize}
  \item BME280 temperature: 0.01$^\circ$C resolution
  \item Arduino ADC: 5V / 1024 = 4.88mV resolution (10-bit)
  \item Higher resolution = more precise measurements.
\end{itemize}

\section{SAMPLING RATE}

\begin{itemize}
  \item Frequency of sensor measurements. Our system: 1 Hz (one reading per second).
  \item Higher sampling rate = more data, faster battery drain, larger log files.
\end{itemize}

SENSOR FUSION (Already Defined)\par

\begin{itemize}
  \item See Section 3 - Combining multiple sensor inputs for better accuracy.
\end{itemize}

\section{TIMESTAMP}

\begin{itemize}
  \item Time marker for each data point. Arduino uses millis() function (milliseconds
  \item since power-on). Format: Unix timestamp or formatted date/time string.
\end{itemize}

\section{SECTION 6: SOFTWARE \& FIRMWARE TERMINOLOGY}

\section{ARDUINO IDE}

\begin{itemize}
  \item Software development environment for programming Arduino boards. Uses simplified
  \item C/C++ language. Includes code editor, compiler, and upload tool. Free and
  \item open-source.
\end{itemize}

ArduPilot\par

\begin{itemize}
  \item Open-source autopilot software for fixed-wing, multirotor, and other vehicles.
  \item Features: GPS navigation, autonomous missions, advanced sensors. Alternative
  \item firmware for our MATEK F405-WMN (in addition to iNav).
\end{itemize}

\section{BOOTLOADER}

\begin{itemize}
  \item Small program that allows firmware updates via USB. Pre-installed on Arduino
  \item and flight controllers. Runs first on power-up, checks for upload request.
\end{itemize}

\section{CONFIGURATOR}

\begin{itemize}
  \item GUI software for flight controller setup. Options:
  \item iNav Configurator (for iNav firmware)
  \item Mission Planner (for ArduPilot)
  \item Betaflight Configurator (for Betaflight)
  \item Used to configure PID values, modes, calibration, OSD layout.
\end{itemize}

\section{FIRMWARE}

\begin{itemize}
  \item Low-level software running directly on microcontroller. Read-only memory,
  \item flashed via USB. Examples:
  \item iNav on flight controller
  \item Arduino sketch on Arduino Nano
  \item ELRS on receiver
\end{itemize}

\section{FLASHING}

\subsection{Uploading firmware to microcontroller's flash memory. Process}

\begin{itemize}
  \item 1. Connect device via USB
  \item 2. Select COM port in software
  \item 3. Click "Flash" button
  \item 4. Wait for completion (device reboots automatically)
\end{itemize}

\section{GITHUB}

\begin{itemize}
  \item Online platform for hosting open-source code. ELRS, iNav, ArduPilot all hosted
  \item on GitHub. Can download latest versions, report bugs, contribute code.
\end{itemize}

iNav\par

\begin{itemize}
  \item Open-source flight control firmware optimized for fixed-wing and long-range
  \item flight. Features GPS navigation, missions, return-to-home. Alternative to
  \item ArduPilot, easier setup. Recommended for our project.
\end{itemize}

\section{LIBRARY}

\begin{itemize}
  \item Pre-written code that provides functionality for common tasks. Arduino libraries
  \item for our project:
  \item Adafruit\_BME280: BME280 sensor interface
  \item SD: MicroSD card file operations
  \item Wire: I2C communication
\end{itemize}

\section{MISSION PLANNING}

\subsection{Creating pre-programmed flight path with GPS waypoints. Software}

\begin{itemize}
  \item Mission Planner (ArduPilot)
  \item iNav Configurator (iNav)
  \item Define takeoff, waypoints (lat/long/alt), loiter circles, landing.
\end{itemize}

\section{OSD LAYOUT}

\begin{itemize}
  \item Configuration of on-screen display elements (position, size, visibility).
  \item Can show: battery voltage, current, altitude, GPS position, speed, heading,
  \item artificial horizon, etc.
\end{itemize}

\section{SKETCH}

\begin{itemize}
  \item Arduino program file (.ino extension). Contains setup() function (runs once)
  \item and loop() function (runs repeatedly). Our data logging sketch reads sensors
  \item and writes to SD card.
\end{itemize}

USB BOOTLOADER (See BOOTLOADER)\par

\section{WAYPOINT}

\begin{itemize}
  \item GPS coordinate (latitude, longitude, altitude) defining point in mission.
  \item Aircraft flies from waypoint to waypoint autonomously. Can specify actions
  \item at each waypoint (loiter time, heading, speed).
\end{itemize}

\section{SECTION 7: SAFETY \& TESTING TERMINOLOGY}

\section{BALANCE CHARGING}

\begin{itemize}
  \item LiPo charging method that ensures all cells reach same voltage. Uses balance
  \item connector (JST-XH) to monitor individual cell voltages. REQUIRED for safe
  \item LiPo charging. Never charge without balance lead connected!
\end{itemize}

\section{BENCH TEST}

\subsection{Testing on workbench before installation in aircraft. Advantages}

\begin{itemize}
  \item Safer (no propeller)
  \item Easier troubleshooting (access to all components)
  \item Can measure voltages and currents
  \item Always bench test before flight test!
\end{itemize}

\section{CELL VOLTAGE}

\subsection{Voltage of individual cell in multi-cell battery. LiPo safe voltages}

\begin{itemize}
  \item Full charge: 4.2V per cell
  \item Nominal: 3.7V per cell
  \item Safe minimum: 3.3V per cell
  \item DANGER: 3.0V per cell (damages cell permanently)
  \item 4S battery: 16.8V full, 14.8V nominal, 13.2V minimum
\end{itemize}

\section{CONTINUITY TEST}

\begin{itemize}
  \item Using multimeter to verify electrical connection (wire not broken). Set meter
  \item to $\Omega$ mode, touch probes to both ends of wire. Should beep/read \textasciitilde{}0$\Omega$ if continuous.
\end{itemize}

\section{FAILSAFE TEST}

\subsection{Verifying aircraft behavior when radio link is lost. Process}

\begin{itemize}
  \item 1. Arm aircraft on ground
  \item 2. Turn off transmitter
  \item 3. Verify failsafe activates (motor cuts, servos go to preset positions)
  \item 4. Turn transmitter back on
  \item 5. Verify control restored
\end{itemize}

\section{FIRE EXTINGUISHER}

\begin{itemize}
  \item Required safety equipment for LiPo batteries. Class D (metal fire) or ABC
  \item extinguisher. LiPo fires extremely hot and dangerous. Alternative: Sand bucket
  \item or metal container to smother fire.
\end{itemize}

FUSE\par

\begin{itemize}
  \item Protective device that breaks circuit if current exceeds rated value. Should
  \item be installed on battery positive lead. Our system: 40A fuse protects against
  \item short circuits, prevents battery damage and fire risk.
\end{itemize}

\section{GROUND TEST}

\subsection{Final system check before flight}

\begin{itemize}
  \item Power on sequence (battery last)
  \item GPS lock verification
  \item Radio range test (walk 50m away, verify control)
  \item Failsafe test
  \item Control surface deflection check (full travel, correct direction)
\end{itemize}

\section{LIPO-SAFE BAG}

\begin{itemize}
  \item Fire-resistant bag for charging and storing LiPo batteries. Contains fire if
  \item battery catches fire during charging. NOT fireproof (will eventually burn through),
  \item but provides time to move outside. Alternative: Metal ammo can.
\end{itemize}

\section{PUFFED BATTERY}

\begin{itemize}
  \item Damaged LiPo battery with swollen cells (bulging). Caused by overcharge, over-
  \item discharge, physical damage, or age. [WARNING] EXTREMELY DANGEROUS - can explode or
  \item catch fire! Dispose immediately at battery recycling center.
\end{itemize}

\section{RANGE TEST}

\begin{itemize}
  \item Verifying radio control range before flight. ELRS: Built-in range test mode
  \item reduces TX power to simulate long distance. Walk away from aircraft until warning
  \item (should be >100m). Verify control and failsafe work correctly.
\end{itemize}

\section{SMOKE TEST}

\begin{itemize}
  \item Informal term: First power-on of new build. "If it doesn't smoke, it works!"
\end{itemize}

\subsection{Serious first power-up procedure}

\begin{itemize}
  \item 1. Visual inspection (no shorts, correct polarity)
  \item 2. Connect battery slowly, watch for sparks
  \item 3. Feel components for excessive heat
  \item 4. Measure voltages with multimeter
  \item 5. If anything smokes/smells: DISCONNECT IMMEDIATELY!
\end{itemize}

\section{STORAGE VOLTAGE}

\begin{itemize}
  \item Safe voltage for long-term LiPo storage: 3.8V per cell (15.2V for 4S).
  \item Reduces stress on cells, prolongs battery life. Use charger's "storage mode"
  \item if not flying for >1 week.
\end{itemize}

\section{THERMAL RUNWAY}

\begin{itemize}
  \item Uncontrolled exothermic reaction in damaged LiPo battery. Cell overheats, causing
  \item chemical breakdown, generating more heat, causing neighboring cells to fail.
  \item Results in fire/explosion. Prevention: Proper charging, damage inspection,
  \item voltage monitoring.
\end{itemize}

\section{VOLTAGE ALARM}

\begin{itemize}
  \item Small device that beeps when LiPo voltage drops below set threshold. Plugs into
  \item balance connector. Prevents over-discharge during flight. Set alarm at 3.5V/cell
  \item (14.0V for 4S) for safety margin.
\end{itemize}

\section{SECTION 8: CONNECTORS \& MECHANICAL TERMINOLOGY}

\section{BULLET CONNECTOR}

\subsection{Cylindrical push-fit electrical connector. Common sizes}

\begin{itemize}
  \item 2mm: Low current signal wires
  \item 3.5mm: Motor phase wires (our ESC to motor)
  \item 4mm: High current battery connections
  \item Male = pin (on wire), Female = socket (on wire or component)
\end{itemize}

\section{CRIMPING}

\begin{itemize}
  \item Mechanical joining of connector to wire using special tool. Alternative to
  \item soldering for some connectors (JST, Dupont). Requires proper crimp tool for
  \item connector size.
\end{itemize}

\section{DUPONT CONNECTOR}

\begin{itemize}
  \item 0.1" pitch plastic housing with crimped pins. Used for breadboard connections,
  \item Arduino headers, sensor wiring. Can make custom cable lengths. Separate male
  \item and female pins.
\end{itemize}

HEAT SHRINK (Already Defined)\par

\begin{itemize}
  \item See Section 1 - Insulating tubing that shrinks when heated.
\end{itemize}

JST CONNECTOR (Already Defined)\par

\begin{itemize}
  \item See Section 2 - Small standardized connector family.
\end{itemize}

\section{MOUNTING TAPE / FOAM TAPE}

\begin{itemize}
  \item Double-sided foam tape for securing electronics. Provides vibration dampening
  \item and electrical insulation. Use for mounting FC, GPS, receiver.
\end{itemize}

\section{SERVO ARM / SERVO HORN}

\begin{itemize}
  \item Plastic attachment on servo output shaft. Connects to pushrod/linkage for
  \item control surface actuation. Usually cross-shaped with multiple holes for different
  \item throw distances.
\end{itemize}

\section{SERVO EXTENSION}

\begin{itemize}
  \item Cable to extend servo wire length. Standard 3-pin, male-female. Available in
  \item various lengths (6", 12", 24"). Allows mounting servos far from FC.
\end{itemize}

\section{SOLDER / SOLDERING}

\subsection{Joining metal parts with molten metal alloy (tin-lead or lead-free). Process}

\begin{itemize}
  \item 1. Heat joint with soldering iron (300-350$^\circ$C)
  \item 2. Apply solder (melts and flows into joint)
  \item 3. Remove iron, let cool (2-3 seconds)
  \item Good solder joint: Shiny, smooth, no cold spots. Bad joint: Dull, cracked,
  \item blobby (redo it).
\end{itemize}

\section{TWIST TIE / ZIP TIE / CABLE TIE}

\begin{itemize}
  \item Plastic strap for bundling and securing wires. Prevents wires from snagging
  \item on moving parts, keeps electronics neat. Use locking zip ties (not reusable)
  \item for permanent installation.
\end{itemize}

\section{WIRE STRIPPING}

\begin{itemize}
  \item Removing insulation from wire end to expose conductor. Use wire strippers
  \item (adjustable for gauge) or knife (carefully). Strip enough for connection but
  \item not excessive (keeps joint mechanically strong).
\end{itemize}

XT30 / XT60 (XT60 Already Defined)\par

\subsection{High-current connector series}

\begin{itemize}
  \item XT30: 30A rated, smaller/lighter (good for UBEC power input)
  \item XT60: 60A rated, standard for 3S-6S batteries (our main battery connector)
  \item XT90: 90A rated, heavy duty (overkill for our project)
\end{itemize}

\section{SECTION 9: TROUBLESHOOTING TERMS}

\section{COLD SOLDER JOINT}

\begin{itemize}
  \item Poor solder connection that looks dull/grainy instead of shiny. Caused by moving
  \item joint before solder cools, or insufficient heat. Results in intermittent connection.
  \item Fix: Reheat with soldering iron until shiny and smooth.
\end{itemize}

\section{GROUND LOOP}

\begin{itemize}
  \item Unwanted current flow between multiple ground paths at different voltages.
  \item Causes noise, erratic behavior, flickering OSD. Prevention: Star grounding
  \item topology (all grounds to ONE point).
\end{itemize}

\section{INTERMITTENT CONNECTION}

\subsection{Electrical connection that works sometimes but not always. Causes}

\begin{itemize}
  \item Cold solder joint
  \item Loose connector
  \item Broken wire strand (partially broken)
  \item Corrosion
  \item Troubleshooting: Wiggle wires while powered to reproduce problem.
\end{itemize}

\section{MAGIC SMOKE}

\begin{itemize}
  \item Humorous term for electronics component failure. "I let the magic smoke out =
  \item I broke it." Expensive lesson - always double-check polarity and connections!
\end{itemize}

\section{NO GPS FIX / NO GPS LOCK}

\begin{itemize}
  \item GPS module cannot determine position. Troubleshooting:
  \item Verify clear sky view (no roof, no dense trees)
  \item Check antenna not damaged/disconnected
  \item Verify TX/RX wires not swapped
  \item Wait longer (cold start can take 5+ minutes)
  \item Check baud rate matches FC configuration
\end{itemize}

\section{OSCILLATION}

\begin{itemize}
  \item Unstable rapid shaking or vibration. In flight: PID values too high (reduce P
  \item and D gains). In servo: Mechanical binding or damaged feedback potentiometer.
\end{itemize}

\section{POWER CYCLING}

\begin{itemize}
  \item Turning device off and on again. Basic troubleshooting step that clears transient
  \item errors, resets firmware. For our system: Disconnect battery, wait 5 seconds,
  \item reconnect.
\end{itemize}

\section{REVERSED SERVO}

\begin{itemize}
  \item Servo moves opposite direction from command (up command = down movement). Fix
  \item in FC configurator: Enable "reverse" checkbox for that channel. Do NOT physically
  \item reverse servo wires (changes null position).
\end{itemize}

\section{RX LOSS / LINK LOST}

\subsection{Radio receiver not receiving transmitter signal. Causes}

\begin{itemize}
  \item Transmitter off or too far away
  \item Receiver antenna damaged
  \item Frequency interference
  \item Low transmission power
  \item Result: Failsafe activates. Check RSSI (signal strength) telemetry.
\end{itemize}

SHORT CIRCUIT (Already Defined)\par

\begin{itemize}
  \item See Section 1 - Direct connection between + and -, bypassing load.
\end{itemize}

\section{VOLTAGE SPIKE}

\begin{itemize}
  \item Brief surge in voltage above normal level. Caused by inductive kickback (motor
  \item stopping), switching transients. Can damage sensitive electronics. Mitigation:
  \item Capacitors across power supply, TVS diodes on signal lines.
\end{itemize}

\section{SECTION 10: ACRONYMS QUICK REFERENCE}

ADC     Analog-to-Digital Converter\par

AHRS    Attitude and Heading Reference System\par

ARM     Advanced RISC Machine (processor architecture)\par

AWG     American Wire Gauge\par

BEC     Battery Eliminator Circuit\par

BLDC    Brushless DC Motor\par

CRSF    Crossfire Protocol\par

CSV     Comma-Separated Values\par

DC      Direct Current (as opposed to AC - Alternating Current)\par

DOF     Degrees of Freedom\par

ELRS    ExpressLRS (radio control protocol)\par

EMI     Electromagnetic Interference\par

ESC     Electronic Speed Controller\par

FC      Flight Controller\par

FPV     First Person View\par

GND     Ground (0V reference)\par

GPS     Global Positioning System\par

GUI     Graphical User Interface\par

I2C     Inter-Integrated Circuit\par

IDE     Integrated Development Environment\par

IMU     Inertial Measurement Unit\par

JST     Japan Solderless Terminal (connector type)\par

KV      RPM per Volt constant (motor rating)\par

LED     Light Emitting Diode\par

LiPo    Lithium Polymer (battery type)\par

mAh     Milliamp-hours (battery capacity)\par

MCU     Microcontroller Unit\par

NMEA    National Marine Electronics Association\par

OSD     On-Screen Display\par

PCB     Printed Circuit Board\par

PID     Proportional-Integral-Derivative (control algorithm)\par

PPM     Pulse Position Modulation\par

PWM     Pulse Width Modulation\par

RC      Radio Control\par

RF      Radio Frequency\par

RSSI    Received Signal Strength Indicator\par

RTH     Return to Home\par

RX      Receiver (or Receive pin on serial)\par

SBUS    Serial Bus (FrSky protocol)\par

SCL     Serial Clock (I2C)\par

SDA     Serial Data (I2C)\par

SPI     Serial Peripheral Interface\par

TX      Transmitter (or Transmit pin on serial)\par

UAV     Unmanned Aerial Vehicle\par

UART    Universal Asynchronous Receiver-Transmitter\par

UBEC    Universal Battery Eliminator Circuit\par

UBX     u-blox binary GPS protocol\par

USB     Universal Serial Bus\par

VOC     Volatile Organic Compounds (what MQ135 detects)\par

\section{APPENDIX: COMMON FORMULAS}

\section{OHM'S LAW:}

\begin{itemize}
  \item V = I $\times$ R    (Voltage = Current $\times$ Resistance)
  \item I = V / R    (Current = Voltage / Resistance)
  \item R = V / I    (Resistance = Voltage / Current)
\end{itemize}

\section{POWER:}

\begin{itemize}
  \item P = V $\times$ I         (Power = Voltage $\times$ Current)
  \item P = I² $\times$ R        (Power = Current squared $\times$ Resistance)
  \item P = V² / R        (Power = Voltage squared / Resistance)
\end{itemize}

\section{BATTERY CAPACITY:}

\begin{itemize}
  \item Runtime (hours) = Capacity (mAh) / Current (mA)
  \item Energy (Wh) = Capacity (Ah) $\times$ Voltage (V)
\end{itemize}

WIRE CURRENT CAPACITY (Approximate):\par

\section{12 AWG: 20-25A}

\section{14 AWG: 15-20A}

\section{16 AWG: 10-15A}

\section{18 AWG: 7-10A}

\section{20 AWG: 5-7A}

\section{22 AWG: 3-5A}

\section{24 AWG: 2-3A}

\begin{itemize}
  \item 26 AWG: 1-2A (signal wires only)
\end{itemize}

\section{SERVO TRAVEL:}

\begin{itemize}
  \item Travel (degrees) = (Pulse Width - 1.5ms) $\times$ 90$^\circ$ per 0.5ms
  \item Center: 1.5ms = 0$^\circ$
  \item Min: 1.0ms = -90$^\circ$
  \item Max: 2.0ms = +90$^\circ$
\end{itemize}

\section{LIPO CELL VOLTAGES:}

\begin{itemize}
  \item Full charge:    4.20V per cell (16.8V for 4S)
  \item Nominal:        3.70V per cell (14.8V for 4S)
  \item Storage:        3.80V per cell (15.2V for 4S)
  \item Low warning:    3.50V per cell (14.0V for 4S)
  \item Critical min:   3.30V per cell (13.2V for 4S)
  \item Damaged:        <3.00V per cell (DO NOT USE!)
\end{itemize}

\section{END OF TERMINOLOGY REFERENCE}

This document should be used as a quick reference guide throughout the project.\par

When you encounter an unfamiliar term in meeting materials, wiring diagrams, or\par

troubleshooting, refer to this glossary.\par

Recommended: Print this document and keep it accessible during build sessions.\par

\subsection{For questions not answered here, consult}

\begin{itemize}
  \item Arduino reference: arduino.cc
  \item iNav documentation: github.com/iNavFlight/inav/wiki
  \item ELRS documentation: expresslrs.org
  \item Component datasheets (available on manufacturer websites)
\end{itemize}

Good luck with your build!\par

\end{document}
